\documentclass[OPS,lsstdraft,authoryear,toc]{lsstdoc}
% GENERATED FILE -- edit this in the Makefile
\newcommand{\lsstDocType}{RTN}
\newcommand{\lsstDocNum}{107}
\newcommand{\vcsRevision}{491e809-dirty}
\newcommand{\vcsDate}{2025-09-05}


% ---------------------------------------------------------------------------
% DRAFT: technote-style build of the Rubin ToO MDC paper. The ApJS (AASTeX)
% version of this manuscript lives in Rubin-ToO-dev/paper/main.tex; this file
% is the same text converted to the lsstdoc class. Red text still needs input.
% ---------------------------------------------------------------------------
\usepackage{booktabs}
\newcommand{\tbd}[1]{\textcolor{red}{[#1]}}
\newcommand{\todo}[1]{\textcolor{red}{\textbf{[TODO: #1]}}}
\newcommand{\ToO}{ToO\xspace}
\newcommand{\mdc}{MDC\xspace}

\title{An End-to-End Mock Data Challenge for Target-of-Opportunity Response with the Vera C. Rubin Observatory: Gravitational-Wave and Neutrino Alerts}
% TODO: working title; confirm

\hypersetup{
    pdftitle={An End-to-End Mock Data Challenge for Target-of-Opportunity Response with the Vera C. Rubin Observatory},
    pdfauthor={Sean Patrick MacBride},
    pdfkeywords={gravitational waves, neutrinos, target of opportunity, Rubin Observatory, LSST, mock data challenge}
}

% TODO: ORCID, full author list and order, Rubin / LVK / IceCube publication-policy review
\author{Sean P. MacBride}

\setDocRef{RTN-107}
\setDocUpstreamLocation{\url{https://github.com/lsst/rtn-107}}

\date{\vcsDate}

\setDocAbstract{%
The Vera C.\ Rubin Observatory will devote about 3\% of the ten-year LSST to
target-of-opportunity (ToO) observations, including follow-up of gravitational-wave (GW) events
and high-energy neutrinos. We present the Rubin ToO Mock Data Challenge (\mdc),
an end-to-end exercise in which \tbd{number} mock alerts from external experiments were
delivered to the Rubin ToO system and the resulting observing response was simulated with the
Rubin scheduler. Combining the simulated observations with simulated electromagnetic (EM)
counterpart light curves generated with LightCurveLynx, we measure
\tbd{headline detection fraction for GW MDC v1}, \tbd{headline for GW MDC v2}, and
\tbd{headline for neutrino alerts}, and characterise the system latency between alert trigger and
receipt on the summit (\tbd{median latency}). We discuss the implications for the
LSST ToO strategy and for the design of future challenges.
}

\setDocChangeRecord{%
  \addtohist{1}{2026-10-23}{Draft skeleton, converted from the ApJS (AASTeX) manuscript.}{MacBride}
}

\begin{document}

\maketitle

\section{Introduction}\label{sec:intro}

The discovery of the kilonova AT\,2017gfo following GW170817 \citep{Abbott2017MMA} demonstrated the
scientific value of rapid, wide-field optical follow-up of multi-messenger sources. The Vera C.\ Rubin
Observatory \citep{Ivezic2019} combines a $\sim$9.6\,deg$^2$ field of view with a large collecting area,
and was designed with a dedicated ToO program of roughly 3\% of survey time
\citep{Andreoni2022,MacBride2026}. ToO observations sit outside the normal survey mode of the LSST,
so their performance depends on a chain of components that must work together: alert
ingestion, conversion of an external localisation into a scheduler-readable footprint, the
Rubin scheduler's response, and the on-sky observations themselves.

\todo{1--2 paragraphs: Rubin ToO program history (Andreoni et al.\ 2022 ToO white paper / ToO 2024 workshop recommendations); state of the system in the first year of operations (MacBride et al.\ 2026, arXiv:2607.00217).}

Before and during early operations, such a chain is best validated end to end with a mock data
challenge. In a \mdc, external experiments inject simulated alerts into the same interfaces that real
alerts will use, and the downstream response is measured against known truth.
\todo{Describe the MDC concept, who ran it (LVK, IceCube, Rubin), dates, and how this relates to RTN-107.}

This paper reports the results of the Rubin ToO \mdc{} for GW and neutrino alerts. Our contributions are:
\begin{enumerate}
  \item a description of the \mdc{} alert streams (two versions of the LVK \mdc{} and a neutrino simulation);
  \item a reproducible pipeline that turns each alert into a simulated Rubin observing response using the Rubin scheduler;
  \item a definition of EM counterpart ``detection'' that combines the simulated observations with simulated light curves; and
  \item measured detection fractions, response times, and system latencies for each alert stream, with a discussion of implications for the LSST ToO program.
\end{enumerate}

The paper is organised as follows. Section~\ref{sec:data} describes the data; Section~\ref{sec:methods} the
methods, including our definition of a detection; Section~\ref{sec:results} the results; and
Sections~\ref{sec:discussion} and~\ref{sec:conclusion} the discussion and conclusions.

%% ---------------------------------------------------------------------------
\section{Data}\label{sec:data}

Table~\ref{tab:data} summarises the data products used in this work.

\begin{table}[ht]
\centering\small
\caption{Data products used in this work.\label{tab:data}}
\begin{tabularx}{\textwidth}{l l X l}
\toprule
Product & Source & Content & Used in\\
\midrule
LVK \mdc{} v1 alerts & LVK \tbd{pipeline / production} & \tbd{N} events; localisation skymaps & \S\ref{sec:res_v1}\\
LVK \mdc{} v2 alerts & LVK \tbd{pipeline / production} & \tbd{N} events; localisation skymaps & \S\ref{sec:res_v2}\\
Neutrino alerts & IceCube \tbd{stream} & \tbd{N} events; localisations & \S\ref{sec:res_nu}\\
EFD alert records & Rubin EFD & trigger and receipt timestamps; footprint & \S\ref{sec:latency}\\
Simulated observations & Rubin scheduler & per-visit opsim tables & all\\
Simulated counterparts & LightCurveLynx & light curves per injection & all\\
\bottomrule
\end{tabularx}
\end{table}

\subsection{Gravitational-wave alerts}\label{sec:data_gw}

\todo{Describe MDC v1 and MDC v2: what differs (e.g.\ population, sky-map pipeline, alert content, latency model, case labels). Event counts by source class (BNS, NSBH, BBH, \dots) and by false-alarm-rate threshold.}

For each GW alert the Rubin Engineering Facilities Database (EFD) record contains the trigger
timestamp, the time the alert was received on the summit, an alert type, and a binary HEALPix
footprint (\texttt{NSIDE}=32, nested ordering) of the region the scheduler is asked to cover. The
source of each \mdc{} event is matched to a table of injection properties from the
external experiment, which supplies the true sky position of the injected source
(\texttt{declination}, \texttt{corrected\_ra}) used later to score detections.

Beyond the \mdc{} events themselves, we also use LVK-format \texttt{bayestar} all-sky skymaps from simulated
O4-like and O5-like observing-run populations (O4 HLV, O5a, O5b, O5c) to place the \mdc{}
results in a population context \todo{decide whether this belongs in the paper or in a companion paper}.

\subsection{Neutrino alerts}\label{sec:data_nu}

\todo{Describe the IceCube neutrino alert simulation: alert stream (Gold/Bronze), event rate, localisation (error-ellipse $\to$ footprint), how mock alerts were generated and delivered, number of alerts, whether the sky coverage assumes a single event or a multiplet.}

\subsection{Simulated EM observations}\label{sec:data_obs}

The simulated observations are produced with the Rubin scheduler and the \texttt{lsst\_survey\_sim}
framework, using the operational scheduler configuration (\texttt{ts\_config\_scheduler}, the
\texttt{fbs\_config\_lsst\_survey.py} survey configuration and the block-407 deep-drilling-field
generator), together with \texttt{rubin\_nights} utilities for day-obs and observatory
modelling \todo{versions, commits, and citations}. Previously executed visits are loaded from
an \texttt{opsim} table so that the scheduler starts each simulation in the state of the real survey on the
night before the alert. The science programs considered are \texttt{BLOCK-365}, \texttt{BLOCK-407}
and \texttt{BLOCK-408}.

\subsection{Simulated EM counterparts}\label{sec:data_lc}

\todo{LightCurveLynx: package version and citation; source models used per alert class (e.g.\ kilonova grid for BNS/NSBH; neutrino counterparts such as supernovae/TDEs/blazar flares); priors on peak luminosity, redshift/distance (drawn from the GW distance posterior where applicable), extinction; number of injections per event; time sampling.}

\section{Methods}\label{sec:methods}

\subsection{From alert to scheduler footprint}\label{sec:footprint}

For simulated skymaps, the 90\% credible region is computed from the probability density with the
greedy credible-level algorithm \citep[\texttt{ligo.skymap};][]{Singer2016}. The resulting
mask is down-sampled to \texttt{NSIDE}=32, and any coarse pixel containing any part of
the 90\% region is retained, so that the footprint is binary and inclusive. This reproduces the
operation performed on the summit. For \mdc{} events the footprint is taken
directly from the EFD record. A pointing centre for the ToO is drawn at random from within
the footprint \todo{confirm: in current scripts, the centre is a randomly selected footprint pixel; state whether this affects the response}.

\subsection{Scheduler simulation}\label{sec:sim}

Each alert is wrapped in a \texttt{TargetoO} object (trigger time, duration, footprint, centre and
ToO type) and passed to the scheduler through a simulated ToO server. One ToO is simulated
per run. The simulation starts on the day-obs preceding the trigger and runs for a number of nights that
depends on the alert type: 20 nights for GW events, 50 for binary black holes, 30 for neutrinos, and 5
each for solar-system and Galactic-supernova alerts. For skymap-driven simulations of
O4/O5-like events a fixed duration of 20 nights is used \todo{one script uses 10; confirm which applies where}.
Multi-night runs include modelled downtime and clouds; single-night runs do not.
\todo{State seeing model (default vs.\ fixed), whether real EFD downtime is replayed, and the band scheduler configuration.}

\subsection{Observing strategies}\label{sec:strategies}

Two ToO strategies, labelled ``gold'' and ``silver,'' are mapped to the scheduler ToO types
\texttt{GW\_case\_B} and \texttt{GW\_case\_D}, respectively. \todo{Define each case in physical terms: bands, number of visits per field, cadence over the follow-up, depth, area-based triggers.}
For the population simulations, trigger times are drawn uniformly over one year (MJD 61041--61407) to average over
the observing-season and lunar-phase dependence of the response.

\subsection{Counterpart light curves}\label{sec:lc_methods}

\todo{How LightCurveLynx is configured and evaluated against the simulated visits: observing-band bookkeeping, per-visit 5$\sigma$ depth (from the opsim $m_5$ column), photometric noise model, treatment of non-detections, host/extinction, and any peculiarities of sampling light curves on irregular, multi-band visit tables.}

\subsection{Definition of a detection}\label{sec:detection}

\todo{Fix and state the detection criterion. Candidate formulation to confirm or replace:}
an injected counterpart is \emph{covered} if its true position falls inside a simulated visit footprint,
and \emph{detected} if, among covered visits, the simulated flux exceeds the visit's 5$\sigma$ limiting
magnitude on at least $N_{\rm det}=\tbd{N}$ epochs in at least $N_{\rm band}=\tbd{N}$ bands within
$\Delta t=\tbd{window}$ of the trigger. We report separately
(i) the \emph{sky-coverage fraction}, (ii) the \emph{coverage-and-depth fraction}, and (iii) the \emph{detection fraction}
after the full criterion, and we state the dependence on each threshold in Section~\ref{sec:discussion}.
\todo{Also: how detection is scored against GW localisation (probability-weighted coverage) vs.\ against injected position.}

\subsection{System latency}\label{sec:latency}

For each \mdc{} alert we compute the interval between the event trigger timestamp and the time at
which the alert was recorded in the EFD on the summit, and report its distribution.
\todo{Decompose into external-experiment latency and Rubin-internal latency if the timestamps allow; relation to the Rubin system latency analysis.}

\subsection{Reproducibility and computing}\label{sec:compute}

Each simulation is an independent single-core job (16\,GB memory, 30\,min wall-time request) run as a SLURM
array over the list of alerts, with one array task per skymap. The simulation scripts are available at
\url{https://github.com/seanmacb/Rubin-ToO-dev} \todo{archive with DOI (Zenodo) and freeze versions}.

\section{Results}\label{sec:results}

\subsection{LVK MDC version 1}\label{sec:res_v1}
\todo{Event table; sky-coverage, coverage-and-depth and detection fractions vs.\ source class, distance, and localisation area; time-to-first-observation distribution; latency; example sky map with executed visits (figure).}

\subsection{LVK MDC version 2}\label{sec:res_v2}
\todo{Same quantities as v1; direct comparison of v1 and v2, and attribution of any differences (population, strategy, system changes).}

\subsection{Neutrino simulations}\label{sec:res_nu}
\todo{Coverage and detection fractions for the neutrino footprints; sensitivity to the footprint size and the assumed counterpart class; time-to-first-observation.}

\subsection{Population-level estimates}\label{sec:res_pop}
\todo{Optional: O4-like and O5a/b/c gold vs.\ silver strategy comparison from the population simulations. Remove if out of scope.}

\section{Discussion}\label{sec:discussion}
\todo{(1) What the challenge showed works and what failed; (2) sensitivity of the detection fraction to the definition thresholds and light-curve models; (3) limitations: one ToO per simulation, no competing ToOs or follow-up of earlier alerts, simplified weather/downtime, no image-subtraction or transient-vetting efficiency, light-curve model uncertainty; (4) latency budget relative to kilonova fading timescales; (5) recommendations for ToO strategy and for MDC v3.}

\section{Conclusion}\label{sec:conclusion}
\todo{Short summary of headline numbers and recommendations; data and code availability.}

\appendix

\section{Acknowledgements}

This material is based upon work supported in part by the National Science Foundation through Cooperative Agreements AST-1258333 and AST-2241526 and Cooperative Support Agreements AST-1202910 and AST-2211468 managed by the Association of Universities for Research in Astronomy (AURA), and the Department of Energy under Contract No.\ DE-AC02-76SF00515 with the SLAC National Accelerator Laboratory managed by Stanford University.
Additional Rubin Observatory funding comes from private donations, grants to universities, and in-kind support from LSST-DA Institutional Members.

% TODO: LVK and IceCube acknowledgement wording; funding; computing resources (University of Zurich science cluster).

\section{Facilities and software}

\paragraph{Facilities.} Rubin:Simonyi (LSSTCam), LIGO, Virgo, KAGRA, IceCube

\paragraph{Software.} rubin\_scheduler, rubin\_sim, lsst\_survey\_sim, rubin\_nights, LightCurveLynx, healpy \citep{Zonca2019}, ligo.skymap \citep{Singer2016}, astropy, numpy, pandas, matplotlib

% TODO: add astropy/numpy/matplotlib citations per AAS software policy

\section{References} \label{sec:bib}
\renewcommand{\refname}{} % Suppress default Bibliography section
\bibliography{references,lsst,lsst-dm,refs_ads,refs,books}

\end{document}
